% !TEX root = /home/eb/.vscode/extensions/qftrules.programme-colle/tmp/Exercice.tex
%%%%%%%%%%%%%%%%%%%%%%%%%%%%%%%%%%%
\begin{exo}[PCSI][1][colle]{Rotation d'un œuf dur}
	\Notions{\item Solide indéformable \task Énerge mécanique}
Lorsqu'on impulse une rotation rapide à un œuf
dur posé sur une surface plane, au bout de
quelques tours, l'œuf se dresse et se met à tourner sur sa pointe ! Lorsqu'il perd peu à
peu de la vitesse, il finit par se remettre en position couchée.
%, position où son centre de
% gravité est le plus bas.
% \begin{fig}[]
% 	\centering
% 	\includegraphics[width=0.4\textwidth]{œuf_2.jpg}
% 	\capt{}
% \end{fig}
On souhaite établir pour l'œuf dur la condition de basculement de la position horizontale à la position verticale. On adopte le paramétrage
% $de la figure \ref{fig:œuf} $
ci-dessous :
\begin{fig}[]
	\centering
	\includegraphics[width=0.7\columnwidth]{oeuf_1.png}
	% \capt{\label{fig:œuf} Passage de la position horizontale (à gauche) à la position verticale (à droite)}
\end{fig}
% On ne considère que les états initial et final sans s'intéresser au mécanisme transitoire du redressement de l'œuf.
On modélise l'œuf dur par un ellipsoïde de révolution homogène de masse $m$, de demi petit axe $a$ et de
demi grand axe $b>a$. Le centre de masse $G$ est au centre de l'ellipsoïde. Les moments d'inertie d'un ellipsoïde de masse $m$ par rapport à son axe de rotation $Oz$ s'écrivent :
% \begin{itemize}
% \item
 $J_H=\textstyle\frac{1}{5}m(a^2+b^2)$ lorsque l'œuf tourne à l'horizontal,
% \item
 $J_V=\textstyle\frac{2}{5}ma^2$ lorsque l'œuf tourne à la verticale.
% \end{itemize}
On pose $\Omega$ la vitesse angulaire de rotation de l'œuf.

\begin{quest}
%question ajoutée par rapport au sujet
\QR{Comparer les deux moments d'inertie $J_H$ et $J_V$ et commenter physiquement.}
{Comme $b>a$, on a $J_H>J_V$. En effet, la masse est globalement répartie plus loin de l'axe de rotation, donc le moment d'inertie est plus grand.}

\QR{
Exprimer l'énergie mécanique totale de l'œuf dans les deux positions $E_{m_H}$ et $E_{m_V}$. On choisira comme origine de l'énergie potentielle celle d'altitude nulle.
}{
Dans les deux cas, $E_m=E_c+E_p=\frac{1}{2}J\Omega^2+mgz_G$, on a donc
\eq{
\boxed{E_{m_H}=\frac{1}{2}J_H\Omega^2+mga} \text{ et }\boxed{E_{m_V}=\frac{1}{2}J_V\Omega^2+mgb}
}
}

\QR{ \label{q:pulsLimite}
Montrer qu'il existe une pulsation limite $\Omega_c$ telle que pour $\Omega>\Omega_c$, la position verticale est d'énergie inférieure à la position horizontale et assure le basculement d'une position à l'autre. On exprimera $\Omega_c$ en fonction de $a$, $b$, et $g$.
}{ \label{q:pulsLimite}
On cherche la condition sur $\Omega$ pour avoir $E_{m_V}<E_{m_H}$, ce qui équivaut à :
\eq{
\frac{1}{2}J_V\Omega^2+mgb<\frac{1}{2}J_H\Omega^2+mga \Rightarrow 2mg (b-a)<(J_H-J_V)\Omega^2 \Rightarrow 2mg (b-a)<\frac{m}{5}\underbrace{(b^2-a^2)}_{(a-b)(a+b)}\Omega^2,
}
ce qui donne 
$$
\Omega >\sqrt{\frac{10 g}{a+b}}.
$$
On a donc $\Omega >\Omega_c$ avec 
\eqbox{
	\Omega_c=\sqrt{\dfrac{10 g}{a+b}}
	}
}

\QR{Calculer $\Omega_c$ pour $a=\SI{2,0}{cm}$, $b=\SI{3,0}{cm}$ et $g=\SI{10}{m.s^{-2}}$. Commenter.
% le résultat obtenu en utilisant les descriptions de l'expérience du document.
}{
L'application numérique donne $\Omega_c=\SI{45}{rad.s^{-1}}$, soit environ 7 tours/s. Cette valeur est du bon ordre de grandeur puisqu'on nous parle dans le document d'une vitesse de rotation d'une dizaine de tours par seconde.
}
\end{quest}
\end{exo}
%%%%%%%%%%%%%%%%%%%%%%%%%%%%%%%%%%%


%%%%%%%%%%%%%%%%%%%%%%%%%%%%%%%%%%%
\begin{exo}[PC][1][colle]{Le pendule pesant}
On considère un pendule pesant constitué d’une barre homogène de longueur $\ell$ et de masse $m$ suspendu à un point fixe $O$ par un fil inextensible de masse négligeable. Le pendule évolue dans le champ de pesanteur $\Vec{g}$. On note $\theta$ l'angle que fait le pendule avec la verticale. On suppose que le pendule est soumis à une force de frottement fluide $\Vec{f} = -\lambda \Vec{v}$.
\begin{quest}
	\item Déterminer le moment d’inertie $J$ du pendule par rapport à l’axe de rotation $O$.
	\item Exprimer le moment cinétique $\Vec{L}$ du pendule par rapport à l’axe de rotation $O$, ainsi que les moments des forces par rapport à ce même axe.
	\item Montrer que l'équation différentielle du mouvement s'écrit :
	\[
	\ddot{\theta}+\omega_0^2 \sin\theta+\omega_0 Q \dot{\theta}= 0,
	\]
	où $\omega_0$ et $Q$ sont à déterminer.
	\item Résoudre cette équation différentielle pour une condition initiale $\theta(0)=\theta_0$ et $\dot{\theta}(0)=0$.
	
\end{quest}
	
\end{exo}
%%%%%%%%%%%%%%%%%%%%%%%%%%%%%%%%%%%
